% Restructured 2026-09-28 from method_tests.tex (verbatim move, then edited).
\subsection{Does the conditioning matter? Only for methods that solve}
\label{subsec:ridge}

\S\ref{subsec:conditioning} found a four-order-of-magnitude difference in
$\kappa(\bar H)$ between the two initialisation regimes. A difference that
size which nothing observable depends on would be a curiosity, so we tested
it, pre-registering the prediction before running the cells.

\paragraph{First, the null, and what it can actually detect.} We compare
every method on matched cohorts, same method, same base, shared against
independent initialisation, with three cohorts per arm so the
pre-registered gate can be applied: a cell counts as an effect only if the
mean difference exceeds $\max(0.005, 2\,\mathrm{SE})$. Across $20$
base-by-method cells, $15$ show no detectable effect, $5$ favour
independent initialisation, and none favours shared.

The number that matters alongside that is the smallest effect the design
could have found, because a null without one is not evidence of absence.
The gate sits at $0.0052$ nats in the median cell and $0.0109$ in the
worst, so the honest statement is \emph{no effect detectable at about
$0.01$ nats}, not \emph{no effect}. It also disposes of a smaller number
that has been quoted in this setting, a mean difference of $0.0014$ nats
against a metric resolution of about $0.001$: that difference sits below
the median gate, so it is a quantity this design cannot resolve in either
direction.

The five cells that do clear the gate are worth naming rather than
averaging away, because four of them are not spread at random. Three are
the subspace-alignment method, on Mistral, Qwen and Yi, at $0.0096$ to
$0.0129$ nats; the others are TIES on Llama at $0.0151$ and TVQ on Yi at
$0.0087$. Task arithmetic and DARE, the two methods closest to a plain
mean, show nothing anywhere. That ordering is the one the mechanism below
predicts: a method that explicitly operates on the subspace geometry is
the heuristic most exposed to a change in that geometry, and it gains from
the cohort whose subspaces are not near-collinear. The effect is small,
consistent in direction, and we do not build anything on it.

\paragraph{The same null in a unit we do not disclaim.} Everything above is
worst-task NLL excess, and \S\ref{sec:discussion} reports that this metric does
not track downstream accuracy across methods and on one benchmark ranks them
backwards. That objection lands harder on a null than on an effect: the
differences here are hundredths of a nat, which is the range in which we have
told the reader not to trust the metric. Meanwhile the positive result beside
it, the solver collapse, has two units. Reporting one half in a disputed unit
and the other half in two is an asymmetry we would rather remove than defend,
so we pre-registered the null in accuracy as well: first at one cohort per
arm, where no cell cleared a binomial gate of about $6$ points on GSM8K and
$11$ on HumanEval (Appendix~\ref{subsec:heur-n1}), and then, because a null at
that resolution is weak, as an equivalence test at three cohorts per arm
(\texttt{prereg\_heuristics\_tost}).

\begin{table}[t]
\centering
\small
\caption{The heuristics in downstream accuracy at three cohorts per arm, as an
equivalence test. $\bar d$ is the mean over the four bases of
$\mathrm{acc}_{\mathrm{indep}} - \mathrm{acc}_{\mathrm{shared}}$, so a negative
value favours the shared arm; the interval is the $90\%$ interval of the two
one-sided tests, and ``margin'' is the smallest equivalence margin the method
passes. The registered margins are $0.03$ on GSM8K and $0.05$ on HumanEval, and
every method passes both. GSM8K exact match at $n = 500$, HumanEval pass@1 at
$n = 164$; $240$ cells.}
\label{tab:heur-tost}
\begin{tabular}{lcccccc}
\toprule
 & \multicolumn{3}{c}{GSM8K} & \multicolumn{3}{c}{HumanEval} \\
method & $\bar d$ & $90\%$ interval & margin & $\bar d$ & $90\%$ interval & margin \\
\midrule
task arithmetic & $-0.003$ & $[-0.009, +0.003]$ & $0.009$ & $-0.019$ & $[-0.030, -0.008]$ & $0.030$ \\
TIES & $-0.013$ & $[-0.020, -0.006]$ & $0.020$ & $-0.007$ & $[-0.021, +0.007]$ & $0.021$ \\
DARE & $-0.007$ & $[-0.015, +0.000]$ & $0.015$ & $-0.016$ & $[-0.021, -0.010]$ & $0.021$ \\
TVQ$_2$ & $+0.004$ & $[-0.004, +0.012]$ & $0.012$ & $+0.004$ & $[-0.013, +0.021]$ & $0.021$ \\
repaired KnOTS & $-0.009$ & $[-0.015, -0.002]$ & $0.015$ & $-0.009$ & $[-0.022, +0.005]$ & $0.022$ \\
\bottomrule
\end{tabular}
\end{table}

\textbf{All five heuristics are equivalent across the two regimes} within the
registered margins, $3$ accuracy points on GSM8K and $5$ on HumanEval, and each
would pass at a margin of $0.9$ to $2.0$ points on GSM8K and $2.1$ to $3.0$ on
HumanEval (Table~\ref{tab:heur-tost}). One of the forty base-by-method cells
clears the per-cell gate: TIES on Llama at HumanEval, where the shared arm
scores $6.5$ points higher. It is the cell that came within $97\%$ of the gate
at one cohort per arm, in the same direction. Three of the four registered
predictions therefore hold, and from here on the null is stated as equivalence
rather than as the absence of a detected effect. It is a statement at these
margins: an effect of one or two points fits inside them.

The fourth prediction failed, and it concerns direction. At one cohort per arm,
$26$ of the $38$ non-zero differences had favoured the shared arm
($p = 0.034$); we had not registered that test, so we did not report the lean
as a finding. We registered the sign test for the three-cohort design and
predicted that the lean would not survive. It survives: $29$ of the $38$
non-zero per-base differences favour the shared arm, an exact sign test gives
$p = 0.0017$, and eight of the ten method-by-benchmark means are negative, TVQ
being the exception on both benchmarks. By the registered branch we report it
in its measured direction: \textbf{on these cohorts the averaging and
sparsifying heuristics score slightly higher when the adapters share their
initialisation}, by at most $1.9$ points on average, well inside the
equivalence margins. It runs opposite to the five NLL cells above, all of which
favoured independent initialisation, so the two units disagree about a
direction while agreeing that nothing is large. It is not the solver mechanism,
since no heuristic enters the ill-conditioned directions on either arm
(Table~\ref{tab:mechanism}), and we have no tested explanation for it. For
practice it means that someone merging with a heuristic loses nothing by a
shared initialisation, and may gain a point.

\paragraph{Then the mechanism, measured.} The explanation for the other
$15$ cells is that those methods never go near the ill-conditioned
direction. That is a measurement rather than an assertion, and it needs no
forward passes: every merge here is weight
arithmetic on the adapter factors. For each merged solution we report its
norm relative to the mean task delta, and the share of its energy lying in
the bottom quartile of $\bar H$'s eigen-directions, in the projected
coordinates of \S\ref{subsec:conditioning}.

\begin{table}[t]
\centering
\small
\caption{What each solution does with the ill-conditioned directions.
Amplification is $\|\Delta_{\mathrm{merged}}\|/\|\bar\Delta\|$; tail
fraction is the share of squared norm in the bottom quartile of $\bar H$'s
spectrum. Ranges are over the four base models. On the independent arm a
tail fraction near $0.21$ is what a generic direction gives when a quarter
of the dimensions holds a quarter of the energy, so the signature is the
$0.73$ to $0.80$ on the shared arm, not the baseline it is measured
against.}
\label{tab:mechanism}
\begin{tabular}{lcccc}
\toprule
 & \multicolumn{2}{c}{shared} & \multicolumn{2}{c}{independent} \\
solution & amplification & tail & amplification & tail \\
\midrule
mean task delta            & 1.00 & 0.000 & 1.00 & 0.21 \\
task arithmetic            & 0.99 & 0.000 & 0.86--0.92 & 0.20 \\
DARE                       & 0.98--0.99 & 0.000 & 0.86--0.90 & 0.20 \\
TIES                       & 1.98--2.04 & 0.000 & 2.04--2.22 & 0.21 \\
TVQ ($b{=}2$)              & 1.95--2.16 & 0.000 & 1.68--1.95 & 0.20 \\
\addlinespace
RegMean, minimal ridge     & 35.4--51.4 & 0.73--0.79 & 3.41--3.67 & 0.29 \\
our encoder, $\lambda = 0$ & 36.7--56.7 & 0.78--0.80 & 3.42--3.67 & 0.29 \\
our encoder, $\lambda^\star$ & 2.20--2.56 & 0.000 & 2.84--3.05 & 0.28 \\
\addlinespace
exact interpolator $\tauH$ & 43.4--64.1 & 0.62--0.66 & 3.96--4.00 & 0.28 \\
\bottomrule
\end{tabular}
\end{table}

Table~\ref{tab:mechanism} gives it. On the ill-conditioned cohort every
heuristic sits at one to two times the mean task delta with a tail
fraction that rounds to zero at three decimals, while both unregularised
solvers sit at $35$ to $57$ times with roughly three quarters of their
energy in the bottom quartile of the spectrum. They are the only methods
that go there, which is why they are the only methods the conditioning
reaches. One incidental check falls out: RegMean at minimal ridge and our
encoder at $\lambda = 0$ agree to two decimals on the independent arm
across all four bases, so two differently derived solvers land on the same
solution once the problem is well posed.

\paragraph{So we sweep the ridge, on two solvers.} To expose the
conditioning one needs a method that actually solves the system. Ours
does, and it carries a Tikhonov ridge, so sweeping $\lambda$ over seven
values, $0$, $0.01$, $0.03$, $0.05$, $0.13$, $0.30$ and $1.00$, on both
cohorts with rank pinned to $16$ in every cell isolates the effect.

\begin{table}[t]
\centering
\small
\caption{Worst-task NLL excess for the unregularised solve
($\lambda = 0$), and the gain from the best ridge. Conditioning is
invisible to the heuristics above and severe here.}
\label{tab:ridge}
\begin{tabular}{lccccc}
\toprule
base & \multicolumn{3}{c}{$\lambda = 0$} &
\multicolumn{2}{c}{ridge gain $L(0) - L(\lambda^\star)$} \\
 & shared & independent & ratio & shared & independent \\
\midrule
Llama-3.1-8B & 0.3081 & 0.0713 & $4.3\times$  & 0.2216 & 0.0000 \\
Mistral-7B   & 9.7799 & 0.1504 & $65\times$   & 9.7301 & 0.0952 \\
Qwen2.5-7B   & 0.1973 & 0.0256 & $7.7\times$  & 0.1867 & 0.0114 \\
Yi-1.5-9B    & 0.2490 & 0.1209 & $2.1\times$  & 0.2211 & 0.0904 \\
\bottomrule
\end{tabular}
\end{table}

Table~\ref{tab:ridge} gives the result. Unregularised, the same merge is
$2.1$ to $65$ times worse on the ill-conditioned cohort; on Mistral it
collapses outright to $9.78$ nats against $0.15$. The ridge gain is larger
on the shared arm on all four bases, which was the pre-registered
prediction and held.

The other half of the prediction failed. We also predicted that $\lambda^\star$ itself would be at least ten times
larger on the shared arm. It is not: $\lambda^\star$ is \emph{identical}
across arms on Mistral, Qwen and Yi ($0.13$, $0.13$, $0.30$), separating
only on Llama. How much regularisation is needed does not track
conditioning in any simple way; how much it buys does. By the
pre-registered rules this makes the test \textsc{partial}, and we claim
only the half that held.

\subsection{Both arms at three cohorts}
\label{subsec:shared-arm}

Table~\ref{tab:ridge} rests on one shared cohort against three independent
ones, so its noise gate borrowed its standard error entirely from the
independent side. A shared-versus-independent contrast measured that way is
three draws against one, and a reader is right to discount it. We registered
the completion in advance and ran the sweep on two further shared cohorts, so
both arms are $n = 3$ and the gate carries a term from each. This is a
replication of a result we had already published, so it could only cost us
something; the registration says so and fixes the consequence text for the
outcome where it fails.

It does not fail. With
$\mathrm{SE} = \sqrt{\mathrm{SE}_s^2 + \mathrm{SE}_i^2}$ and the gate at
$\max(0.005, 2\,\mathrm{SE})$, the ridge gain is larger on the shared arm on
\textbf{four of four} bases and the unregularised excess is worse there by at
least a factor of two on \textbf{four of four}: ratios of $4.42$, $76.2$,
$10.8$ and $2.82$, against gates of $0.024$, $4.26$, $0.128$ and $0.166$.
Mistral's shared arm averages $11.72$ nats at $\lambda = 0$ against $0.154$
independent.

One quantity is worth reporting because it could have embarrassed the earlier
table and does not. The published single-cohort value, \texttt{seed1}, sits
$0.23$, $0.53$, $0.75$ and $0.66$ standard deviations from its own arm's
three-cohort mean. We registered that a distance above $2$ on two or more
bases would oblige us to say the published value was unrepresentative of its
own arm. It is representative, and Table~\ref{tab:ridge}'s numbers stand as
single-cohort estimates of quantities the three-cohort means confirm.

\subsection{A second adapter-built solver: RegMean without data}
\label{subsec:regmean}

Everything above concerns one solver, and we built it. The first question a
referee asks is whether the effect belongs to the family or to our
implementation. RegMean~\citep{jin2023regmean} solves a least-squares system in
per-task input-activation Grams; its \emph{data-free, adapter-only} form, which
replaces each activation Gram by $A_t^\top A_t$, is a second solver whose
matrix is built from the adapters' own factors. We say plainly that this form
is a variant we constructed, not the published method, and it is exactly the
kind of solver the mechanism says is exposed; \S\ref{subsec:data-solvers}
tests the published method with real activation Grams, and two further
methods we did not build.

We swept it on the same grid and cohorts, registering two predictions: the
ridge gain larger on the shared arm, and the minimally regularised excess worse
there by at least a factor of two, each on at least three of four bases. The
numbers below use the Gram the method is defined with; an earlier version of
this table was computed with a fallback Gram by mistake, and
Appendix~\ref{subsec:gram-error} gives the account and both sets of numbers.

RegMean ships a regulariser of its own, and it is \emph{not} the Tikhonov
ridge swept here: it scales the non-diagonal entries of each Gram by
$\alpha$, default $0.9$. The $\lambda I$ in this table is ours, and
\S\ref{subsec:regmean-default} evaluates RegMean's own default end to end.

\begin{table}[t]
\centering
\small
\caption{Adapter-only RegMean on both cohorts, with the described Gram
$\sum_t A_t^\top A_t$. The minimally regularised reference is
$\lambda = 10^{-6}$ rather than $0$: the Gram has rank at most $Tr = 64$ in
$d = 4096$ dimensions, so $\lambda = 0$ is singular rather than
ill-conditioned. Our encoder solves inside $\range(\bar H)$, so its
$\lambda = 0$ is defined; each method is compared only against itself across
arms.}
\label{tab:regmean}
\begin{tabular}{llcccc}
\toprule
base & cohort & $\lambda = 10^{-6}$ & $\lambda^\star$ & gain & ratio \\
\midrule
Llama-3.1-8B & shared      & 0.2918 & 0.03 & 0.1781 & \multirow{2}{*}{$4.1\times$} \\
             & independent & 0.0710 & $10^{-6}$ & 0.0000 & \\
\addlinespace
Mistral-7B   & shared      & 10.1651 & 0.05 & 10.1155 & \multirow{2}{*}{$65.6\times$} \\
             & independent & 0.1551 & 0.05 & 0.1006 & \\
\addlinespace
Qwen2.5-7B   & shared      & 0.1504 & 0.05 & 0.1391 & \multirow{2}{*}{$5.8\times$} \\
             & independent & 0.0258 & 0.13 & 0.0117 & \\
\addlinespace
Yi-1.5-9B    & shared      & 0.2179 & 0.13 & 0.1947 & \multirow{2}{*}{$1.8\times$} \\
             & independent & 0.1206 & 0.13 & 0.0879 & \\
\bottomrule
\end{tabular}
\end{table}

Both predictions hold: the ridge gain is larger on the shared arm on four of
four bases, and the unregularised excess is at least twice as large there on
three of four (Yi reaches $1.8$). The effect is not a property of our encoder.

Two details are worth reporting. The selected ridge is the same on both arms
of Mistral and of Yi, so on this solver too the amount of regularisation
needed does not track the conditioning even though the amount it buys does,
the split our own sweep found. And Llama's independent arm selects the minimal
$\lambda$ and gains exactly zero: on a well-conditioned cohort there is nothing
to repair.

\subsection{Solvers built from data, whitening, and a diagonal}
\label{subsec:data-solvers}

The mechanism makes a sharper prediction than ``solvers are exposed''. The
collapse makes the matrix $\sum_t A_t^\top A_t$ ill-conditioned; a solver is
exposed if it inverts a matrix built from the adapters' own factors, and should
not be if its matrix comes from data or is diagonal, since a diagonal cannot see
overlap between subspaces. We registered that prediction for four published
methods, none of ours, before running any of them
(\texttt{prereg\_regmean\_real}, \texttt{prereg\_third\_solver}): RegMean with
real input-activation Grams at its published $\alpha = 0.9$; MaTS~\citep{tam2024mats}
with its blockwise (KFAC) Fisher solved by 100 conjugate-gradient iterations,
as in the official code; TSV-M~\citep{gargiulo2025tsv}, which whitens the
stacked singular vectors of the task updates and was therefore predicted to be
exposed; and diagonal Fisher merging~\citep{matena2022fisher}. Each ran on all
three cohorts of both arms on all four bases.

\begin{table}[t]
\centering
\small
\caption{Methods we did not build. Worst-task NLL excess, shared over
independent (mean over three cohorts per arm), and whether the difference
clears the registered gate $\max(0.005, 2\,\mathrm{SE})$. The prediction column
is the mechanism's, registered before any cell ran.}
\label{tab:data-solvers}
\begin{tabular}{llcccc}
\toprule
method & predicted & Llama & Mistral & Qwen & Yi \\
\midrule
RegMean, activation Grams & immune & $1.03$ & $1.14$ & $0.99$ & $1.13$ \\
MaTS (KFAC Fisher, CG)    & immune & $0.98$ & $0.94$ & $0.93$ & $0.99$ \\
diagonal Fisher           & immune & $1.08^{\ast}$ & $1.02$ & $0.96$ & $0.98$ \\
TSV-M (whitening)         & exposed & $1.08$ & $1.19^{\ast}$ & $0.81$ & $0.91^{\ast}$ \\
\bottomrule
\multicolumn{6}{l}{\footnotesize $^{\ast}$ clears the gate; for Yi the
difference favours the shared arm.}
\end{tabular}
\end{table}

The three immunity predictions hold (Table~\ref{tab:data-solvers}): RegMean
with activation Grams shows no effect on any base, in NLL or in accuracy
(Table~\ref{tab:regmean-acc}), MaTS none, and the diagonal
Fisher merge one small effect on Llama ($1.08$ times), below the registered
bound. The exposure prediction for TSV-M fails: it shows an effect on Mistral
only ($1.19$ times), the opposite sign on Yi, and nothing on Llama and Qwen, and
in accuracy no cell clears the gate on either benchmark. By the registered
branch, the whitening case was predicted and did not hold, and the exposure
rule is stated for solvers alone.

One registered detail also failed. We predicted that RegMean with activation
Grams would keep the pre-truncation norm of its solution within three times the
mean task delta; it reaches $2.2$ to $6.5$. The amplification is the same on
both arms (for example $3.15$ shared against $3.44$ independent on Llama),
because it comes from the activation Grams themselves, whose condition numbers
reach $10^{4}$ to $10^{7}$ whatever the initialisation. The immunity to
initialisation holds; the bound we put on the norm was wrong.

\paragraph{The rule this licenses.} Across eleven methods, nine of them not ours,
the collapse reaches exactly the merges that solve a linear system in a matrix
built from the adapters' own factors: our encoder and adapter-only RegMean.
It does not reach averaging or sparsifying heuristics, which are equivalent
across regimes within a few accuracy points (\S\ref{subsec:ridge}),
solvers whose matrix comes from data, a diagonal solver, or whitening.

\subsection{Does RegMean's published default already fix this?}
\label{subsec:regmean-default}

Our practical recommendation reduces to regularising a solve-based merge, and
RegMean ships a regulariser. If its default were already strong enough on a
collapsed cohort, the recommendation would shrink to the much smaller claim
that the audit tells you when the default matters. The question is worth
answering rather than leaving for a referee, and it is a lookup plus a
CPU-only computation.

RegMean's published regulariser scales the non-diagonal entries of each Gram
by $\alpha$, default $0.9$, which is equivalent to adding
$(1-\alpha)\operatorname{diag}(G)$: a shrinkage whose size the data sets,
rather than a free parameter. It does regularise, and on a rank-deficient Gram
it does make the system solvable. The question is whether it is enough.

On the shared-initialisation cohorts it is not, by about an order of
magnitude. Averaged over sampled layers, $\alpha = 0.9$ is equivalent to an
isotropic ridge of $5.2\times10^{-4}$ to $6.1\times10^{-4}$ across the four
bases, against the $\lambda^\star$ of $0.03$ to $0.13$ that the sweep selects on
the shared arm (Table~\ref{tab:regmean}): $50$ to $250$ times smaller. The
conditioning it leaves behind reflects that. The $\alpha$-shrunk system has condition number $6.6\times10^{3}$ to
$9.1\times10^{3}$, while $\lambda = 0.01$ brings the same layers to about
$1.5\times10^{2}$, a factor of $55$ apart. The default is on the wrong side of
the gap this paper is about.

The independent arm is the control, and it behaves as the story requires. The
\emph{same} default leaves $\kappa$ at $1.2\times10^{3}$ to $1.4\times10^{3}$
there, six to seven times lower than on the shared arm, while the equivalent
ridge it applies is identical to three significant figures
($5.2$ to $6.1\times10^{-4}$ on both arms). That is the point in one line:
$\alpha$-shrinkage is a fixed dose set by the data's scale and not by the
cohort's conditioning, so it neither knows nor cares which regime it is in.
On a well-conditioned cohort the dose is ample. On a collapsed one it is
far too small, and nothing in the method notices.

\paragraph{End to end.} The conditioning argument was a computation, so we
registered the merge itself (\texttt{prereg\_regmean\_real}): adapter-only
RegMean with $\alpha = 0.9$ and no ridge, on all three cohorts of both arms, in
NLL and in downstream accuracy. In NLL both registered predictions hold on three
of four bases. The shipped default is worse on the shared arm by $18$, $68$ and
$5.3$ times on Llama, Mistral and Yi, each past the gate; on Qwen it is $41$
times worse on average but one of the three shared cohorts sits far from the
other two ($2.54$ nats against $0.20$ and $0.22$), which widens the gate beyond
the difference. The audit's starting ridge, $\lambda = 0.01$, improves on the
default on the shared arm past the gate on the same three bases.

\begin{table}[t]
\centering
\small
\caption{RegMean end to end in downstream accuracy, mean over three cohorts per
arm. The first two column pairs are the adapter-only form (Gram
$\sum_t A_t^\top A_t$) at the shipped $\alpha = 0.9$ and at the audit's starting
ridge $\lambda = 0.01$; the last is the published method with activation Grams
at $\alpha = 0.9$. GSM8K exact match at $n = 500$, HumanEval pass@1 at
$n = 164$; $144$ cells.}
\label{tab:regmean-acc}
\begin{tabular}{llcccccc}
\toprule
 & & \multicolumn{2}{c}{adapter-only, $\alpha{=}0.9$}
   & \multicolumn{2}{c}{adapter-only, $\lambda{=}0.01$}
   & \multicolumn{2}{c}{activation Grams} \\
base & benchmark & shared & indep. & shared & indep. & shared & indep. \\
\midrule
Llama-3.1-8B & GSM8K     & 0.168 & 0.705 & 0.477 & 0.730 & 0.784 & 0.792 \\
             & HumanEval & 0.187 & 0.541 & 0.425 & 0.553 & 0.646 & 0.659 \\
Mistral-7B   & GSM8K     & \textbf{0.000} & 0.285 & 0.029 & 0.323 & 0.532 & 0.519 \\
             & HumanEval & \textbf{0.000} & 0.057 & 0.002 & 0.093 & 0.360 & 0.366 \\
Qwen2.5-7B   & GSM8K     & 0.148 & 0.611 & 0.580 & 0.625 & 0.624 & 0.617 \\
             & HumanEval & 0.510 & 0.683 & 0.589 & 0.685 & 0.726 & 0.728 \\
Yi-1.5-9B    & GSM8K     & 0.093 & 0.724 & 0.628 & 0.726 & 0.773 & 0.775 \\
             & HumanEval & 0.077 & 0.130 & 0.154 & 0.130 & 0.354 & 0.348 \\
\bottomrule
\end{tabular}
\end{table}

In accuracy the picture is sharper (Table~\ref{tab:regmean-acc}). On the shared
arm the shipped default loses most of GSM8K on every base, scoring $0.09$ to
$0.17$ where the same method reaches $0.61$ to $0.72$ on the independent arm,
and on Mistral it scores \textbf{zero on both benchmarks in all three shared
cohorts}. The starting ridge beats the default on the shared arm past the gate
on three of four bases on GSM8K (by $31$, $43$ and $54$ points on Llama, Qwen
and Yi) and on two on HumanEval (Llama and Qwen), so the registered accuracy
prediction holds. Mistral is the exception and we name it: $\lambda = 0.01$
does not rescue it ($0.029$ and $0.002$), and Table~\ref{tab:regmean} selects
$\lambda^\star = 0.05$ for this solver on that base, so the starting ridge is
where a sweep begins rather than an answer. The published method with
activation Grams, the control, scores the same on both arms in all eight
base-by-benchmark cells, within $1.3$ points everywhere, and is the strongest
merge in the table on every base: a practitioner who has the data RegMean was
designed for is unaffected by the collapse.

$\alpha$ is a documented knob, and a practitioner who lowers it is doing what
\S\ref{subsec:ridge} recommends. The claim is only that the shipped default does
not do it for them, and that nothing in the method signals when it should.

\subsection{The same test without rank truncation}
\label{subsec:untruncated}

\S\ref{subsec:conditioning} notes that our pipeline truncates every merge
back to rank $r$, so the sweep above tests conditioning through a
rank-constrained solver while the interpolator whose norm we report has
rank up to $Tr$. Repeating the sweep with truncation disabled removes that
mismatch, and the truncated sweep is then the comparison arm rather than
being re-run.

\begin{table}[htb]
\centering
\small
\caption{The conditioning link with rank truncation disabled. Worst-task NLL
excess in nats at $\lambda = 0$, and the shared-to-independent ratio with and
without truncation. The truncated column is the sweep of
Table~\ref{tab:ridge}, reproduced here as the comparison arm rather than
re-run. \textbf{Two entries here invite a second look and we have taken it.}
Mistral's shared value, $11.6283$, sits exactly $1.0000$ above RegMean's
$10.6283$ in Table~\ref{tab:regmean} at four decimals; the underlying numbers
are $11.628281$ and $10.628309$, so the gap is $0.999973$ and the agreement is
a coincidence between two different methods on two different runs rather than
a transcription. And Qwen's independent value matches
Table~\ref{tab:ridge}'s to four decimals, $0.025571$ against $0.025638$;
that one is not coincidence but the point of the table, since on a
well-conditioned cohort there is almost nothing in the discarded directions
for truncation to remove.}
\label{tab:untruncated}
\begin{tabular}{lcccc}
\toprule
 & \multicolumn{2}{c}{$\lambda = 0$ excess} & \multicolumn{2}{c}{shared / independent} \\
base & shared & independent & untruncated & truncated \\
\midrule
Llama-3.1-8B & 7.4266  & 0.0635 & $116.9\times$ & $4.3\times$ \\
Mistral-7B   & 11.6283 & 0.1582 & $73.5\times$  & $65\times$ \\
Qwen2.5-7B   & 0.8993  & 0.0256 & $35.2\times$  & $7.7\times$ \\
Yi-1.5-9B    & 3.1294  & 0.1227 & $25.5\times$  & $2.1\times$ \\
\bottomrule
\end{tabular}
\end{table}

Table~\ref{tab:untruncated} shows the link surviving on all four bases, with the effect \emph{larger}
without truncation, by factors of $25$ to $117$ against $2.1$ to $65$.
Truncation was not manufacturing the result; it was hiding part of it,
which is what one should expect when all but the leading $16$ directions
of a rank-$64$ interpolator are discarded and the ill-conditioned subspace
goes with them. The numbers we report as primary are therefore the
conservative ones.

\paragraph{What this licenses, and what it does not.} It licenses a
conditional recommendation: if you merge with a method that solves a
linear system, check your cohort's conditioning and regularise, because
unregularised it can fail by a factor of $65$ at rank $16$ and $117$
without truncation, on two solvers, one of which we did not build. It does
\emph{not} license the general advice to initialise independently. On the
methods most practitioners actually use the null above stands, in NLL and
as equivalence in accuracy, with two small qualifications that point in
opposite directions: in NLL the subspace-alignment method gains a little
from independent initialisation on three of four bases, and in accuracy
the heuristics as a group lean slightly toward the shared arm. We do not
dress either up as anything more.

\subsection{The collapse is visible in downstream accuracy}
\label{subsec:downstream-cond}

Everything above is measured in worst-task NLL excess, and
\S\ref{sec:discussion} reports that this metric does not predict downstream
accuracy across merge methods, and on HumanEval predicts it backwards. A
reader is entitled to extend that objection to the whole paper: if the metric
does not track model quality, why should any of these numbers matter?

The objection is about differences of hundredths of a nat between methods. The
conditioning effect is not that. At $\lambda = 0$ the shared-versus-independent
gap runs from $0.22$ to $11.56$ nats across the four bases at $n = 3$ per
arm ($0.13$ to $9.63$ from the single-cohort Table~\ref{tab:ridge}, which is
the same quantity measured with one draw), one to three orders of magnitude above
both the metric's resolution and the differences the correlation was computed
over. An effect that size either shows up in accuracy or the metric is
measuring something other than model quality. We registered the test without
knowing which, and fixed in advance that a null would move this limitation into
\S\ref{sec:intro} and delete the practice recommendations rather than defend
them.

\begin{table}[t]
\centering
\small
\caption{Downstream accuracy of the same merges, with the repaired scorers.
GSM8K exact match at $n = 500$, HumanEval pass@1 at $n = 164$, its whole test
split. Gate is $\max(5\text{ points}, 2\,\mathrm{SE})$ on the binomial. No cell
had a scorer discard: every generation was scored.}
\label{tab:downstream-cond}
\begin{tabular}{lcccccc}
\toprule
& \multicolumn{3}{c}{shared} & \multicolumn{3}{c}{independent} \\
base & $\lambda{=}0$ & $\lambda^\star$ & gain
     & $\lambda{=}0$ & $\lambda^\star$ & gain \\
\midrule
\multicolumn{7}{l}{\emph{GSM8K, exact match}} \\
Llama-3.1-8B & 0.318 & 0.654 & $+0.336$ & 0.708 & 0.712 & $+0.004$ \\
Mistral-7B   & \textbf{0.000} & 0.462 & $+0.462$ & 0.276 & 0.424 & $+0.148$ \\
Qwen2.5-7B   & 0.496 & 0.808 & $+0.312$ & 0.786 & 0.786 & $+0.000$ \\
Yi-1.5-9B    & 0.346 & 0.802 & $+0.456$ & 0.698 & 0.794 & $+0.096$ \\
\addlinespace
\multicolumn{7}{l}{\emph{HumanEval, pass@1}} \\
Llama-3.1-8B & 0.250 & 0.573 & $+0.323$ & 0.518 & 0.530 & $+0.012$ \\
Mistral-7B   & \textbf{0.000} & 0.244 & $+0.244$ & 0.067 & 0.189 & $+0.122$ \\
Qwen2.5-7B   & 0.530 & 0.744 & $+0.213$ & 0.726 & 0.738 & $+0.012$ \\
Yi-1.5-9B    & 0.024 & 0.091 & $+0.067$ & 0.049 & 0.067 & $+0.018$ \\
\bottomrule
\end{tabular}
\end{table}

\paragraph{It shows up, on both benchmarks.} The ridge improves the shared arm
past the gate on \textbf{four of four} bases on both benchmarks. The
difference-in-differences, which is the claim that the ridge helps
\emph{because of} conditioning rather than helping everything everywhere,
clears the gate on four of four on GSM8K and three of four on HumanEval; the
exception is Yi, where both arms sit near the floor of a benchmark it barely
solves at all ($0.024$ to $0.091$). By the registered rule, requiring either
prediction on at least three of four bases on at least one benchmark, the
outcome is \textsc{confirmed}.

\paragraph{The number worth stating plainly.} On Mistral, the shared-cohort
unregularised merge scores \textbf{zero} on both benchmarks. It is not
producing malformed output that a scorer discards: all $500$ GSM8K generations
and all $164$ HumanEval generations are non-empty, and none contains an
extractable answer or passes a test. The model generates fluent text and has
lost the capability entirely. One Tikhonov term on the same adapters recovers
$0.462$ and $0.244$. That is what four orders of magnitude of conditioning buys
in a unit a practitioner cares about.

\paragraph{What this does not license.} It does not show that worst-task NLL
excess predicts accuracy in general, and we do not claim it. The correlation
that fails in \S\ref{sec:discussion} is across merge \emph{methods} at fixed
cohort, and nothing here tests that. What it shows is that the specific effect
this paper is about is not an artifact of the metric, which is the objection it
was registered to answer.


\subsection{The same results with a working translation adapter}
\label{subsec:translation-fix}

Every cohort above contains a translation adapter that did not learn its task
(\S\ref{subsec:scope}). We retrained it with the one-line loader fix, a
shuffled stream, on Llama-3.1-8B and Mistral-7B in the shared and first
independent cohorts, with each cohort's own seeds, and re-ran the core cells on
the resulting cohorts (\texttt{prereg\_new\_cohorts}, Part T).

The fix helps but does not pass its registered test. The retrained adapter now
improves on the base model's translation loss (by $12\%$ on Llama, where the
original was $1.6\%$ worse, and by $6\%$ on Mistral), but another adapter in the
cohort is still better on translation, so by the registered criterion it has
not learned its task and the caveat of \S\ref{subsec:scope} stands as written.

On these cohorts every other registered prediction holds. The regimes are
unchanged (median cosine $0.995$ and $0.994$ shared, $0.047$ independent);
unregularised adapter-only RegMean is $6.0$ and $64.6$ times worse on the shared
arm with the larger ridge gain there; the five heuristics move by less than
$0.02$ nats in all ten cells; and in accuracy the encoder's $\lambda^\star$
recovers the shared arm past the gate on both bases and both benchmarks. On
Mistral the unregularised shared merge again scores $0.000$ on GSM8K and on
HumanEval, and the ridge again recovers $0.468$ and $0.244$. The results of this
section therefore do not depend on the adapter that failed to train, although
the cohort is still not four well-trained adapters.

\subsection{Setting the ridge}
\label{subsec:lambda-rule}

The practical advice is to regularise an adapter-built solve on a collapsed
cohort. How much is the natural next question, and the first-order theory
suggests a data-free answer: cap the condition number of the adapter Gram. We
registered such a rule before evaluating it anywhere
(\texttt{prereg\_lambda\_rule}): per layer, the smallest isotropic ridge that
brings the condition number of $\sum_t A_t^\top A_t$ on its range down to
$k = 100$, a round number fixed in advance, evaluated only on cohorts that no
ridge had been selected on.

The rule fails its main test. On all six collapsed cohorts of
\S\ref{subsec:translation-fix} and \S\ref{subsec:tr-sweep} it is worse than the
fixed $\lambda = 0.01$ (by $0.08$ to $4.2$ nats), although it beats RegMean's
shipped default on five of six. On the separated cohorts it is within $0.005$
nats of the unregularised solve on four of six; in the other two it is better by
more than that, which the registered test also counts as a miss. By the
registered branch, the recommendation stays as it was: on a collapsed cohort,
start near $\lambda = 0.01$ for adapter-only RegMean and sweep, since the
selected ridge varies by base (Table~\ref{tab:regmean}) and does not track the
conditioning. A data-free rule that sets the ridge well remains open.
